\begin{landscape}
\footnotesize
\begin{longtable}{@{}L{0.6cm}L{7.5cm}L{8.0cm}L{7.5cm}@{}}
\caption{Reviewed literature: objective, method/key findings (with project relevance) and limitations.}\label{tab:lit2}\\
\toprule S.No & Objective / outcome & Method, key findings and relevance & Limitations and gaps\\\midrule\endfirsthead
\toprule S.No & Objective / outcome & Method, key findings and relevance & Limitations and gaps\\\midrule\endhead
\bottomrule\endfoot
1 & Adapts a pretrained vision-language model using local/global alignment and text prompts for remote-sensing segmentation. & Relevant to foundation-model adaptation and class semantics in satellite imagery. Quantitative score: not stated in the abstract. & Natural-image priors and prompt choices may introduce domain/class errors; it is not direct evidence for rural infrastructure planning.  \\
2 & Adapts self-supervised vision foundation models with multiscale cross-attention and hyperbolic pixel-level loss for low-annotation segmentation. & Recent direction relevant to the user's DINOv2/foundation-model segmentation interest and cross-domain transfer concerns. Reported result: mean IoU 67.60\% (iSAID), 55.86\% (LoveDA), 80.07\% (Potsdam) and 96.46\% (Vaihingen). & Very recent results need independent reproduction and target-domain evaluation; published benchmark scores cannot be assumed to transfer to Damarakunta.  \\
3 & Combines local sparse attention and compressed global sparse attention with coarse-to-fine feature enhancement. & Relevant to road structures that require both local detail and long-range continuity. Quantitative score: not stated in the abstract. & Results on benchmark datasets need not transfer to Damarakunta; graph topology must be evaluated separately.  \\
4 & Uses a compact global-context block to improve road continuity while limiting computational complexity. & Directly relevant to connected road extraction and resource-efficient inference. Quantitative score: abstract not retrievable from the DOI registries. & Stronger raster connectivity is not equivalent to a verified, traversable GIS network.  \\
5 & Uses multiscale uncertainty consistency and cross teacher-student attention to exploit unlabeled imagery. & Relevant to uncertainty-aware segmentation and reducing the cost of dense annotation. Reported result (open-access preprint version): mean IoU of 74.62\% (ISPRS-Potsdam) and 50.97\% (LoveDA) with 5\% labels, against 72.01\% and 44.20\% without the proposed modules. & Model uncertainty must be calibrated and evaluated; pseudo-label consistency is not a substitute for reference ground truth.  \\
6 & Introduces an open-vocabulary segmentation framework designed for orientation and scale variation in high-resolution remote-sensing images. & Relevant to flexible class definitions and possible future handling of diverse land-cover classes. Quantitative score: abstract not retrievable from the DOI registries. & Open-vocabulary output reliability for small rural objects and project-specific classes requires direct validation.  \\
7 & Plans paths for intelligent vehicles in off-road terrain with several objectives, including terrain-related failure risk and weather. & Actor-critic deep reinforcement learning on a traversability map with meteorology-driven heuristics; reported to beat standard DRL and non-DRL planners in efficiency and convergence stability. Relevant to terrain-aware multi-objective path cost. Quantitative score: not stated in the abstract. & Plans one vehicle path; it does not plan a network over many targets, and it is not an infrastructure-planning method.  \\
8 & Uses directional strip convolutions, multiscale residual features and a combined decoder to reduce fragmented road predictions. & Relevant to road completeness and robustness across image resolutions. Quantitative score: not stated in the abstract. & Improved raster extraction must still be converted to and audited as a connected road graph.  \\
9 & Reviews building extraction across optical and SAR imagery, including footprints, building attributes, height and change-related tasks. & Useful for positioning the building-detection stage and understanding geometric and semantic challenges. Quantitative score: not stated in the abstract. & A review does not establish which model is best for the project's sensor or rural settlement; local benchmarking is required.  \\
10 & Uses neural architecture search with a lightweight backbone and multiscale fusion to balance segmentation accuracy and speed. & Relevant when processing large imagery under limited compute and memory. Reported result: 54.5\% mIoU (LoveDA), 87.8\% (Potsdam) and 84.1\% (Vaihingen) at 132.7 FPS; 72.6\% mIoU at 164.0 FPS (CityScapes) and 72.3\% mIoU at 186.4 FPS (CamVid). & Speed/accuracy trade-offs and LoveDA results are not evidence of Damarakunta performance or downstream planning correctness.  \\
11 & Jointly extracts buildings and roads using cross-task and cross-scale feature interactions. & Highly relevant to the project's coupled building/road layers and the spatial relationship between settlement objects and roads. Reported result (open-access preprint version): F1 higher by 3.21\% (buildings) and 3.38\% (roads) than the baseline (81.85\% and 71.03\%). & Joint extraction still does not solve route optimization, target prioritization, legal access or graph connectivity verification.  \\
12 & Introduces regular and stochastic occlusion-aware decoder components for road regions obscured by dense occlusions. & Addresses a common source of road discontinuity in overhead imagery. Quantitative score: not stated in the abstract. & Occlusion-aware segmentation may still hallucinate or miss segments; independent topology and reference checks remain essential.  \\
13 & Uses optical-image boundary information in a teacher-student distillation design for SAR building extraction. & Shows a way to exploit complementary modalities to improve boundaries where SAR interpretation is difficult. Reported result: F1 higher by 4.3 to 7.2 points and IoU higher by 4.9 to 8.0 points than the strong baseline on MSAW. & Requires suitable paired or aligned modalities and cannot be assumed applicable to an optical-only workflow.  \\
14 & Uses image enhancement and a building-integrity branch to improve footprints under partial occlusion. & Relevant to buildings partly hidden by trees, shadows or other structures. Reported result: IoU +1.49\%, +1.16\%, +2.26\% and F1 +0.82\%, +0.70\%, +1.22\% over HD-Net on WHU, Inria and AIRS. & Benchmark gains do not establish local accuracy; ambiguous merged footprints still need review and error tracking.  \\
15 & Uses building-aware feature clustering and activation refinement to generate better pseudo-masks from weaker labels. & Relevant to reducing labeling effort when preparing rural reference data. Reported result: IoU of 0.8556, 0.8163 and 0.7797 on the three datasets. & Weak supervision can produce uncertain pseudo-labels; these cannot be treated as ground truth for C1 evaluation.  \\
16 & Combines text-guided feature activation and Segment Anything Model refinement for weakly supervised remote-sensing segmentation. & Relevant to reducing dense annotation needs and exploring foundation-model assistance. Quantitative score: not stated in the abstract. & The method is weakly supervised and task/data dependent; outputs require independent label-based evaluation.  \\
17 & Corrects broken road surfaces and inaccurate topology in road extraction from high-resolution imagery. & Three stages: multi-branch road segmentation, energy-based connectivity refinement and topology correction into a vector network; improvements reported on several public datasets. Relevant to vector road graphs and gaps in them. Reported result: F1 improved by 3.3\% to 4.5\% and IoU by 2.0\% to 5.1\% over state-of-the-art methods. & Public benchmark data; whether corrected connections exist physically is not checked, and no route planning is done.  \\
18 & Learns improvement heuristics that approximate the Pareto front for multi-objective route planning problems such as logistics routing. & Population-based deep reinforcement learning model (Pareto Improver); the abstract reports that it outperforms other state-of-the-art methods on several multi-objective routing problems. Shows that multi-objective routing is an active area. Quantitative score: not stated in the abstract. & Vehicle-routing type problems with no terrain cost surface, shared infrastructure or uncertain targets.  \\
19 & Plans an urban power distribution network by combining GIS data with a hybrid optimisation model. & Extracts the street graph from GIS data, computes shortest paths between nodes, then uses Clark and Wright savings plus mixed-integer programming; reports that shortest-path distances cut investment by more than 30\% compared with Euclidean or Manhattan distances in an urban Chinese case. Supports using path distances rather than straight lines in network planning. Reported result: computation more than 1000 times faster than MILP, and an investment decrease of 30\% or more with shortest-path distance instead of Euclidean or Manhattan distance. & Urban power network, not rural roads; inputs are street data, not uncertain remote-sensing targets.  \\
20 & Improves generalisation of land-cover segmentation to unseen regions, sensors and seasons. & Frequency-based optimal style mix (FOSMix) randomises image style in the frequency domain with a consistency regulariser. Relevant to the transfer question behind C3. Quantitative score: not stated in the abstract. & Reports segmentation accuracy under domain shift, not the effect on a downstream plan.  \\
21 & Trains one segmentation model from a single source domain so that it works on unseen target domains. & Category Consistency on Domain Randomization (CCDR) randomises texture and style and adds a consistency loss across source and generated domains. Relevant to the case where no labelled target village is available. Quantitative score: not stated in the abstract. & Segmentation benchmark results only; downstream plan quality is not studied.  \\
22 & Introduces local-consistency/global-diversity regularization for unsupervised domain adaptation of land-cover mapping. & Relevant to the domain-shift problem between benchmark imagery and an unseen rural target village. Quantitative score: not stated in the abstract. & Land-cover adaptation results do not prove building-level or road-planning transfer; target-domain validation remains necessary.  \\
23 & Adds a road-location prediction task and location-guided decoder to improve road continuity. & Directly relevant because road connectivity is a structural requirement, not just a pixel-classification score. Quantitative score: not stated in the abstract. & Extracted connectivity still requires vector graph checks, endpoint snapping and comparison against authoritative roads.  \\
24 & Combines semantic and contour information in an interactive-learning approach for building-footprint extraction. & Supports attention to object contours and the possibility of targeted human review of ambiguous footprints. Quantitative score: not stated in the abstract. & Interactive/manual correction increases annotation effort and needs a documented protocol for reproducibility.  \\
25 & Builds a road network as a graph by predicting road nodes and the connectivity between them instead of skeletonising a segmentation mask. & NodeConnect, trained jointly on nodes and a connectivity map; compared with segmentation methods on DeepGlobe and RoadTracer with pixel and topological precision and recall. Relevant to the topology of the road layer a planner needs. Quantitative score: not stated in the abstract. & Topology is scored on benchmark data; the method does not link houses or facilities to the network.  \\
26 & Extracts roads from satellite images together with partial, incomplete road maps (for example from OpenStreetMap) instead of from scratch. & Two-branch Partial-to-Complete network with a gated self-attention module and a missing-part loss; IoU of 70.71\% on SpaceNet and 75.52\% on OSM data are reported. Close to the project situation of an incomplete TGRAC road layer. Reported result: IoU of 70.71\% (SpaceNet) and 75.52\% (OSM). & Completes road maps; it does not decide how houses should be connected, and rural Indian roads are not evaluated.  \\
27 & Introduces adversarial shape learning and a shape regularizer to improve building footprint geometry and segmentation. & Highlights the importance of object-level geometry metrics in addition to pixel accuracy; relevant to house/building target quality. Reported result (open-access preprint version): mean IoU higher by 2.73\% (Inria) and 3.26\% (MAS) than the baseline, and by 1.51\% over the second-best method on Inria. & Footprint quality still depends on imagery, labels and domain characteristics; it does not evaluate downstream road-planning consequences.  \\
28 & Uses attraction-field representations with CNNs to generate building footprints from remote-sensing images. & Provides a geometry-aware alternative to plain pixel segmentation, relevant to separating and delineating adjacent buildings. Quantitative score: abstract not retrievable from the DOI registries. & Accurate footprint extraction alone does not resolve ambiguous building identity, entrances, legal access or route connectivity.  \\
29 & Uses GIS-aware conditioning for individual-building segmentation in very-high-resolution SAR imagery. & Shows how prior spatial/GIS information can assist individual-object delineation; relevant to combining imagery and geographic context. Reported result (open-access preprint version): F1 of 75.08\% for DeepLabv3-CG on Berlin TerraSAR-X data, F1 and IoU 4.40\% and 5.45\% higher with complete footprints than with sensor-visible segments, and a mean building-height error of 2.39 m. & SAR-specific method and assumptions may not transfer directly to optical imagery or rural Telangana without evaluation.  \\
30 & Combines boundary optimization and multiscale context for building extraction. & Addresses boundary quality and scale variation, both important when converting building predictions to vector footprints. Quantitative score: not stated in the abstract. & Segmentation improvements do not by themselves quantify how missed, merged or false buildings alter route plans.  \\
31 & Models multiple regions and scales to handle buildings with varying sizes and layouts. & Relevant to rural settlements containing small houses and larger public buildings. Quantitative score: not stated in the abstract. & Generalization to different sensors, resolutions and rural morphologies must be measured; no network-planning evaluation is implied.  \\
32 & Combines edge cues with segmentation through a boundary-enhancement module that can be used with encoder-decoder networks. & Supports improving footprint edges across multiple building datasets and architectures. Quantitative score: not stated in the abstract. & Sharper boundaries do not guarantee correct instance separation or a real-world accessible connection.  \\
33 & Combines multitask attention Res-U-Net building-shape extraction with offset vectors and auxiliary POI information. & Illustrates how building geometry and semantic/point-of-interest information can be combined. Reported result: building-roof F1 77.05\% and IoU 63.55\%; whole-building F1 79.02\% and IoU 66.05\%, the best among the baseline models. & Targets high-rise buildings and auxiliary delivery data; transfer to rural houses and public facilities requires caution.  \\
34 & Adapts segmentation models across domains for very-high-resolution remote-sensing images. & Curriculum-style local-to-global adaptation using an entropy-based patch score and semantic- and entropy-level domain classifiers. Relevant to cross-village transfer of segmentation. Reported result (open-access preprint version): best segmentation accuracy of 48.06\% and 52.03\% in the two adaptation directions. & Needs target-domain images and does not evaluate network planning.  \\
35 & Extracts roads from satellite imagery while keeping them connected, because pixel-based methods often fragment roads where buildings or trees occlude them. & Connectivity attention network with strip convolutions and a connectivity attention module; the abstract reports state-of-the-art results on SpaceNet and DeepGlobe. Relevant to producing a connected base road layer. Quantitative score: not stated in the abstract. & Benchmark imagery; connectivity of the output is not tested against a planning network and there is no rural Telangana evidence.  \\
\end{longtable}
\end{landscape}
